\documentclass[11pt,letterpaper]{article}

\usepackage[spanish,es-noshorthands]{babel}
\usepackage{fontspec}
\setmainfont{Source Serif Pro}
\usepackage{microtype}
\usepackage{geometry}
\usepackage{booktabs}
\usepackage{tabularx}
\usepackage{enumitem}
\usepackage{xcolor}
\usepackage{graphicx}
\usepackage{csquotes}
\usepackage{amsmath}
\usepackage{siunitx}
\usepackage[version=4]{mhchem}
\usepackage{tikz}
\usetikzlibrary{arrows.meta,positioning,shapes.geometric,fit,backgrounds,patterns,decorations.pathreplacing}
\usepackage{xurl}
\usepackage[hidelinks]{hyperref}
\usepackage[backend=biber,style=apa,sorting=nyt]{biblatex}
\DeclareLanguageMapping{spanish}{spanish-apa}

\geometry{margin=2.3cm}
% Los rótulos de las figuras usan nodos de ancho fijo con texto alineado, y la
% bibliografía compone en modo sloppy: ambos generan avisos de línea holgada
% que son cosméticos. Se silencian para que `make check` siga siendo sensible
% a los desbordes reales.
\hbadness=10000
\setlength{\emergencystretch}{1em}
\setlength{\parindent}{0pt}
\setlength{\parskip}{0.55em}
\setlist{nosep}
\renewcommand{\topfraction}{0.92}
\renewcommand{\bottomfraction}{0.85}
\renewcommand{\textfraction}{0.06}
\renewcommand{\floatpagefraction}{0.72}
\setcounter{topnumber}{2}
\setcounter{totalnumber}{3}
\addbibresource{referencias.bib}
\AtBeginBibliography{\sloppy}

\sisetup{
  output-decimal-marker = {,},
  group-separator = {.},
  group-minimum-digits = 4,
  detect-all
}

% Colores: una generación de interfaces remotas por color, más el rojo de
% aviso, el verde de la mejora y el gris de los ejes.
\definecolor{cRpc}{RGB}{104,110,126}
\definecolor{cCorba}{RGB}{150,96,52}
\definecolor{cSoap}{RGB}{158,134,38}
\definecolor{cRest}{RGB}{36,102,148}
\definecolor{cGql}{RGB}{182,24,124}
\definecolor{cGris}{RGB}{132,124,112}
\definecolor{cRef}{RGB}{70,90,112}
\definecolor{cAviso}{RGB}{176,42,42}
\definecolor{cVerde}{RGB}{56,120,86}

% --- Figura 1: sucesión de curvas S en carriles ------------------------------
\newcommand{\gx}[1]{((#1)-1980)*0.2385}
% #1 base del carril, #2 color, #3 pendiente k, #4 año del punto de inflexión,
% #5 inicio del dominio, #6 fin del dominio
\newcommand{\curvaS}[6]{%
  \draw[#2, very thick, samples=90, domain=#5:#6, variable=\x]
    plot ({\gx{\x}},{#1+0.95/(1+exp(-(#3)*(\x-(#4))))});%
}
\newcommand{\curvaSp}[6]{%
  \draw[#2, very thick, dashed, samples=40, domain=#5:#6, variable=\x]
    plot ({\gx{\x}},{#1+0.95/(1+exp(-(#3)*(\x-(#4))))});%
}
% #1 base del carril, #2 color, #3 nombre, #4 años de vigencia
\newcommand{\carrilRotulo}[4]{%
  \draw[cGris!35] (0,#1) -- (12.40,#1);%
  \node[left, text width=2.55cm, align=right, font=\scriptsize] at (-0.12,#1+0.46)
       {\textbf{\textcolor{#2}{#3}}\\[-1pt]\textcolor{cGris}{#4}};%
}

% --- Figura 2: cronología y curva de difusión --------------------------------
\newcommand{\tx}[1]{((#1)-2011)*0.820}
% #1 año, #2 altura del tallo, #3 color, #4 texto, #5 ancla del texto
\newcommand{\hito}[5]{%
  \draw[#3, thick] ({\tx{#1}},0) -- ({\tx{#1}},#2);%
  \fill[#3] ({\tx{#1}},0) circle (2.2pt);%
  \node[#5, text width=2.9cm, align=center, font=\scriptsize, inner sep=1.6pt]
       at ({\tx{#1}},#2) {#4};%
}
% #1 y, #2 año final, #3 etiqueta, #4 color
\newcommand{\umbral}[4]{%
  \draw[#4, thick, {|[width=4pt]}-{Stealth[length=4pt]}]
       ({\tx{2012}},#1) -- ({\tx{#2}},#1);%
  \node[right, font=\scriptsize, text=#4, inner sep=2.5pt] at ({\tx{#2}},#1) {#3};%
}
\newcommand{\dx}[1]{((#1)-2014)*0.775}
\newcommand{\dy}[1]{(#1)*0.038}

% --- Figura 3: barras apiladas de tiempo y curva de sensibilidad -------------
% #1 y, #2 ms de ida y vuelta, #3 ms de transferencia, #4 ms de servidor,
% #5 rótulo izquierdo, #6 color del tramo de red, #7 total
\newcommand{\barraT}[7]{%
  \node[left, text width=3.6cm, align=right, font=\scriptsize]
       at (-0.12,#1) {#5};%
  \fill[#6] (0,#1-0.27) rectangle ({(#2)*0.002235},#1+0.27);%
  \fill[#6!35] ({(#2)*0.002235},#1-0.27)
               rectangle ({((#2)+(#3))*0.002235},#1+0.27);%
  \fill[cGris!55] ({((#2)+(#3))*0.002235},#1-0.27)
                  rectangle ({((#2)+(#3)+(#4))*0.002235},#1+0.27);%
  \node[right, font=\scriptsize\bfseries, text=#6]
       at ({((#2)+(#3)+(#4))*0.002235+0.12},#1) {#7};%
}
\newcommand{\px}[1]{(log10(#1)+0.30103)*4.500}
\newcommand{\py}[1]{(#1)*0.300}

% --- Figura 4: N+1 de resolvedores y frente de Pareto ------------------------
\newcommand{\nx}[1]{(#1)*0.062}
\newcommand{\ny}[1]{(#1)*0.042}
\newcommand{\fx}[1]{(#1)*0.062}
\newcommand{\fy}[1]{(#1)*0.045}
% #1 x, #2 y, #3 color, #4 etiqueta, #5 posición
\newcommand{\punto}[5]{%
  \fill[#3] ({\fx{#1}},{\fy{#2}}) circle (2.6pt);%
  \node[#5, font=\scriptsize, text=#3, inner sep=2.8pt, align=center]
       at ({\fx{#1}},{\fy{#2}}) {#4};%
}

\title{\vspace{-1.2cm}\textbf{Cuarenta años pidiendo datos a otra máquina}\\
\large Análisis de la sensibilidad tecnológica de GraphQL como alternativa a
REST, leído sobre la genealogía RPC, CORBA, SOAP y REST}
\author{Carlos Luis Rojas Aragonés\\
\small Gestión del Desarrollo Tecnológico: Estrategias y Análisis de Portfolio}
\date{\small 9 de septiembre de 2026}

\begin{document}

\maketitle
\vspace{-0.6em}

\section{La tecnología que elijo y por qué está en la fase que pide el foro}

Elijo \textbf{GraphQL}, el lenguaje de consulta para interfaces de programación
que Facebook construyó en 2012 y liberó en 2015, y lo analizo como posible
alternativa a REST. Lo elijo por tres razones.

La primera es que está exactamente en la fase que piden los ejemplos del
módulo: una tecnología que ya cruzó el punto de inflexión de su curva S, que
lleva cuatro o cinco años en la parte alta de esa curva y cuya figura de mérito
original está prácticamente agotada, mientras su curva de difusión sigue
subiendo. Esa combinación, rendimiento saturado y adopción creciente, es
justamente donde el análisis de sensibilidad tiene algo que decir.

La segunda es que GraphQL tiene una genealogía documentada de cuarenta años.
No apareció de la nada: es el quinto intento de resolver la misma función, y
los cuatro anteriores fracasaron cada uno contra una restricción activa
distinta. Poder ver cinco curvas S de la misma función, con cuatro asíntotas ya
conocidas, es un lujo que pocas tecnologías ofrecen.

La tercera es personal: decido sobre esto en mi trabajo, y quería comprobar si
la decisión que tomo por intuición aguanta cuando se le declara una figura de
mérito y se le hace la aritmética.

\subsection*{Lo que encontré, dicho por adelantado}

\begin{quote}
\textbf{GraphQL llevó su figura de mérito original hasta el límite teórico en
menos de tres años, y desde entonces el problema que la justificaba se ha ido
disolviendo por otro lado: la red se hizo rápida.} La ventaja de GraphQL sobre
REST, medida en tiempo hasta la vista completa, pasa de un factor 13 en la red
móvil de 2012 a un factor 1,03 entre dos servicios del mismo centro de datos.
Lo que sostiene hoy a GraphQL ya no es el tiempo de respuesta sino el coste de
desarrollo por vista nueva y la gobernanza del esquema. Y sus restricciones
activas se mudaron de sitio: ya no están en la red, están en el servidor.
\end{quote}

\section{La función y la figura de mérito}

Antes de comparar hace falta declarar qué función presta la tecnología y con
qué se mide, que es la disciplina que propone \textcite{deweck2022}: la figura
de mérito es la medida cuantitativa del desempeño en la función prestada,
independiente de la implementación que la entregue.

\textbf{La función es esta:} componer, en un cliente, el conjunto exacto de
datos que necesita una vista, cuando esos datos viven en una o varias máquinas
remotas.

Sobre esa función declaro cuatro figuras de mérito. La primera es la que
gobierna la experiencia y es la que uso como principal:

\begin{equation}
T \;=\; n \cdot \mathrm{RTT} \;+\; \frac{B}{v} \;+\; T_{s}
\label{eq:fom}
\end{equation}

donde $T$ es el \textbf{tiempo hasta la vista completa}, $n$ el número de
viajes de ida y vuelta encadenados, $\mathrm{RTT}$ el tiempo de ida y vuelta de
la red, $B$ los bytes transferidos, $v$ el ancho de banda y $T_{s}$ el tiempo
de trabajo del servidor. Las otras tres son las componentes que un diseño
puede mover:

\begin{itemize}
\item $n$, \textbf{viajes por vista}. Mínimo teórico: 1.
\item $\eta = B_{\text{útil}} / B$, \textbf{eficiencia de carga útil}, la
  fracción de los bytes transferidos que la vista realmente pinta. Máximo
  teórico: 1.
\item $E$, \textbf{esfuerzo de desarrollo por consulta nueva}, en minutos de
  programador. No tiene mínimo teórico, y es la única de las cuatro que no
  depende de la red.
\end{itemize}

Esta separación es la que ordena todo lo que sigue, porque \textbf{cada una de
las cinco generaciones de esta historia optimizó una figura de mérito distinta}
y ninguna de ellas optimizó la de la generación anterior.

\section{Cuatro generaciones antes de GraphQL}

Antes de llegar a GraphQL conviene ver contra qué se detuvo cada una de las
cuatro generaciones anteriores, porque el patrón se repite y es el que anticipa
dónde está hoy el límite de GraphQL. El Cuadro~\ref{cuadro:generaciones} resume
las cinco y la Figura~\ref{fig:curvass} las dibuja como lo que son: cinco
curvas S encadenadas, cada una midiéndose con su propia regla.

\begin{table}[t]
\small
\caption{Cinco generaciones para la misma función. Cada una nació optimizando
una figura de mérito que la anterior había descuidado, y cada una se detuvo
contra una restricción activa distinta.}
\label{cuadro:generaciones}
\begin{tabularx}{\textwidth}{@{}p{2.55cm}p{1.5cm}p{3.5cm}XX@{}}
\toprule
\textbf{Generación} & \textbf{Origen} & \textbf{Figura de mérito dominante} &
\textbf{Restricción que la detuvo} & \textbf{Qué la relevó} \\
\midrule
RPC binario
& 1984
& Transparencia: que la llamada remota se escriba como una local
& Heterogeneidad de máquinas y lenguajes; acoplamiento de la interfaz al
  binario
& Objetos distribuidos \\
\addlinespace[2pt]
CORBA y DCOM
& 1991, 1996
& Interoperabilidad entre lenguajes y proveedores
& Peso del intermediario de objetos y objetos remotos de grano fino, muy
  charlatanes; no atraviesan cortafuegos
& Servicios sobre HTTP \\
\addlinespace[2pt]
SOAP y WS-\hspace{0pt}$\ast$
& 1998, 2003
& Interoperabilidad \emph{entre organizaciones} sobre el puerto 80
& Verbosidad del sobre XML y peso del contrato WSDL y de la pila WS-$\ast$
& Recursos y JSON \\
\addlinespace[2pt]
REST y JSON
& 2000, 2006
& Coste de integración: tiempo hasta la primera llamada correcta
& El recurso tiene forma fija: sobre-entrega y sub-entrega de datos
& Consulta compuesta por el cliente \\
\addlinespace[2pt]
GraphQL
& 2012, 2015
& $\eta$ y $n$: bytes útiles y un solo viaje por vista
& \emph{Activa hoy}: coste de resolución campo a campo en el servidor
& Todavía nada \\
\bottomrule
\end{tabularx}
\end{table}

\subsection*{RPC: la llamada que finge ser local (1984)}

La idea fundacional está en \textcite{birrell1984}: hacer que invocar un
procedimiento en otra máquina se parezca lo más posible a invocarlo en la
propia. La figura de mérito era la transparencia para el programador, y en eso
funcionó. Sun ONC RPC la llevó a producción a finales de los ochenta.

La restricción que la detuvo fue la heterogeneidad. El talón y el esqueleto se
generaban a partir de una descripción de interfaz, el cliente y el servidor
quedaban atados al mismo esquema de serialización, y cualquier cambio en la
interfaz rompía a todos los clientes a la vez. Nadie podía publicar una interfaz
para un tercero desconocido.

\subsection*{CORBA y DCOM: objetos remotos (1991 y 1996)}

CORBA aparece en 1991 desde el Object Management Group y DCOM en 1996 desde
Microsoft \parencite{wikicorba}. La figura de mérito se desplaza a la
interoperabilidad entre lenguajes y proveedores, y técnicamente lo consiguen.

Se detuvieron contra dos cosas. La primera es el peso del intermediario de
objetos: montar un ORB era un proyecto en sí mismo. La segunda, y más letal,
es la que en el módulo llamaríamos una restricción del entorno: los objetos
remotos de grano fino invitan a hacer muchas llamadas pequeñas, y como
$T$ crece linealmente con $n$ según la ecuación~\eqref{eq:fom}, una interfaz
charlatana se vuelve inutilizable en cuanto el RTT deja de ser el de una red
local. Y encima IIOP no pasaba por los cortafuegos corporativos.

\subsection*{SOAP y WS-$\ast$: el puerto 80 como puerta trasera (1998 a 2003)}

XML-RPC en 1998 y luego SOAP resolvieron el problema del cortafuegos por la vía
más pragmática que existía: envolverlo todo en XML y mandarlo por HTTP, que es
el único puerto que ninguna organización cierra \parencite{wikixmlrpc}. SOAP 1.2
llega a recomendación del W3C en junio de 2003 \parencite{w3c2003soap}, con WSDL
para el contrato y una pila de especificaciones WS-$\ast$ alrededor.

La figura de mérito pasó a ser la interoperabilidad entre organizaciones, y esa
sí la resolvió: la banca y los sistemas de reservas siguen funcionando con
SOAP hoy. Pero la restricción activa se volvió el propio sobre. Un mensaje SOAP
con sus cabeceras, sus espacios de nombres y su tipado explícito multiplica por
varias veces la carga útil, y el contrato WSDL hace que publicar un cambio sea
un acto de coordinación entre equipos. En mi tabla de figuras de mérito, SOAP
tiene la $\eta$ más baja de las cinco generaciones.

\subsection*{REST y JSON: el coste de integración por encima de todo (2000 a 2006)}

\textcite{fielding2000} describe REST en su tesis del año 2000, no como un
producto sino como el estilo arquitectónico que el propio HTTP ya encarnaba.
Con JSON como formato, REST optimizó una figura de mérito que las tres
generaciones anteriores habían ignorado: \textbf{el tiempo que tarda un
desarrollador ajeno en hacer su primera llamada correcta.} Sin ORB, sin
generación de código, sin contrato previo, con una URL y un navegador.

Ese fue el salto real y por eso REST ganó. Sigue ganando: en 2025 el 93 por
ciento de los desarrolladores encuestados por Postman declaran usar REST
\parencite{postman2025}.

Y su restricción activa es la que dio origen a GraphQL: \textbf{el recurso
tiene forma fija}. El servidor decide qué campos lleva un usuario, y los lleva
siempre, los use la vista o no. De ahí los dos síntomas clásicos:

\begin{itemize}
\item \textbf{Sobre-entrega}: la vista pide un nombre y una foto, y recibe el
  objeto usuario entero. \textcite{brito2019} migraron siete sistemas reales
  de REST a GraphQL y midieron que la respuesta mediana pasaba de 93,5 campos
  a 5,5. Es decir, \textbf{el 94 por ciento de los campos que REST devolvía no
  se usaban}, y en bytes la reducción mediana fue del 99 por ciento.
\item \textbf{Sub-entrega}: la vista necesita datos anidados y el recurso no
  los trae, así que el cliente encadena peticiones. Eso hace crecer $n$, que en
  la ecuación~\eqref{eq:fom} va multiplicado por el RTT.
\end{itemize}

\begin{figure}[!htb]
\centering
\begin{tikzpicture}

\node[anchor=south west, font=\small\bfseries] at (-2.8,7.05)
     {Cinco curvas S para la misma función, 1980--2032};
\node[anchor=south west, font=\scriptsize\itshape, text=cGris] at (-2.8,6.72)
     {cada carril es el grado de madurez alcanzado por esa generación en su
      \emph{propia} figura de mérito, de 0 a 100 por ciento};

% rejilla de décadas
\foreach \a in {1980,1990,2000,2010,2020,2030} {
  \draw[cGris!28] ({\gx{\a}},-0.30) -- ({\gx{\a}},6.55);
  \node[below, font=\scriptsize, text=cGris] at ({\gx{\a}},-0.30) {\a};
}

% --- Carril 1: RPC binario ---
\carrilRotulo{0}{cRpc}{RPC binario}{Sun ONC RPC, 1984}
\curvaS{0}{cRpc}{0.55}{1990}{1982}{1999}
\node[right, font=\scriptsize\itshape, text=cRpc, text width=4.2cm, align=left]
     at ({\gx{2004}},0.34)
     {techo en 1995: la interfaz queda atada al binario y al lenguaje};

% --- Carril 2: CORBA y DCOM ---
\carrilRotulo{1.30}{cCorba}{CORBA y DCOM}{1991 y 1996}
\curvaS{1.30}{cCorba}{0.52}{1997}{1989}{2006}
\node[right, font=\scriptsize\itshape, text=cCorba, text width=4.2cm, align=left]
     at ({\gx{2009}},1.64)
     {techo en 2002: el ORB pesa y el IIOP no cruza el cortafuegos};

% --- Carril 3: SOAP y WS-* ---
\carrilRotulo{2.60}{cSoap}{SOAP y WS-$\ast$}{1998 y 2003}
\curvaS{2.60}{cSoap}{0.70}{2003}{1997}{2011}
\node[right, font=\scriptsize\itshape, text=cSoap, text width=4.2cm, align=left]
     at ({\gx{2013}},2.94)
     {techo en 2007: el sobre XML pesa más que el dato};

% --- Carril 4: REST y JSON ---
\carrilRotulo{3.90}{cRest}{REST y JSON}{2000 y 2006}
\curvaS{3.90}{cRest}{0.45}{2009}{1999}{2024}
\node[right, font=\scriptsize\itshape, text=cRest, text width=3.5cm, align=left]
     at ({\gx{2021}},4.24)
     {techo en 2016: el recurso tiene forma fija};

% --- Carril 5: GraphQL ---
\carrilRotulo{5.20}{cGql}{GraphQL}{2012 y 2015}
\curvaS{5.20}{cGql}{0.40}{2021}{2012}{2026}
\curvaSp{5.20}{cGql}{0.40}{2021}{2026}{2032}
\fill[cGql] ({\gx{2026}},{5.20+0.95/(1+exp(-0.40*5))}) circle (2.6pt);
\node[left, font=\scriptsize\itshape, text=cGql, text width=4.5cm, align=right]
     at ({\gx{2019.6}},5.93)
     {hoy: 88 por ciento de su techo en $n$ y en $\eta$};

% flechas de relevo: del codo de una generación al arranque de la siguiente
\foreach \p/\q/\a/\b in {1994/1992.8/0/1.30, 2001.2/1999.9/1.30/2.60,
                         2006.1/2004.1/2.60/3.90, 2013.9/2015.5/3.90/5.20} {
  \draw[cGris!75, -{Stealth[length=4.5pt]}, thin]
    ({\gx{\p}},{\a+0.90}) -- ({\gx{\q}},{\b+0.10});
}

\end{tikzpicture}
\caption{La función no cambió en cuarenta años; la figura de mérito, sí. Cada
carril dibuja una logística ajustada a mano sobre las fechas documentadas de
aparición, inflexión y saturación de esa generación, normalizada al techo de su
propia figura de mérito: la transparencia para RPC \parencite{birrell1984}, la
interoperabilidad entre lenguajes para CORBA y DCOM \parencite{wikicorba}, la
interoperabilidad entre organizaciones para SOAP \parencite{w3c2003soap}, el
coste de integración para REST \parencite{fielding2000} y la eficiencia de
entrega para GraphQL \parencite{graphql2015}. Las curvas son elaboración propia
y sirven para leer la forma y el relevo, no para leer un valor de un año
concreto. Las flechas grises marcan el relevo: cada generación arranca cuando
la anterior toca su asíntota.}
\label{fig:curvass}
\end{figure}

\section{Qué trajo GraphQL, y por qué en 2012}

En 2011 y 2012 Facebook estaba reescribiendo sus aplicaciones móviles nativas
después de que su apuesta por la web móvil resultara, en palabras de su propio
fundador, el mayor error estratégico de la compañía. Nick Schrock, Lee Byron,
Dan Schafer y Jonathan Dan construyeron un prototipo llamado SuperGraph que en
el verano de 2012 ya servía el muro de noticias de la aplicación de iOS
\parencite{postman2024historia}.

La idea es una sola y es la que rompe la restricción de REST: \textbf{que sea
el cliente quien declare la forma exacta del resultado}, sobre un esquema
tipado que el servidor publica. Un solo punto de entrada, una sola petición,
una respuesta que tiene la misma forma que la consulta.

El Cuadro~\ref{cuadro:escenario} pone números a lo que eso significa en la
pantalla que motivó el invento.

\begin{table}[t]
\small
\caption{El escenario de cálculo: una pantalla de muro con 20 publicaciones,
su autor y tres comentarios con autor por publicación. Los tamaños de objeto
son mi estimación de un objeto completo típico frente a los campos que la
vista pinta de verdad. Cálculo propio.}
\label{cuadro:escenario}
\begin{tabularx}{\textwidth}{@{}Xrrrr@{}}
\toprule
\textbf{Entidad} & \textbf{Unidades} & \textbf{Objeto} &
\textbf{Campos de} & \textbf{Total} \\
 & & \textbf{completo} & \textbf{la vista} & \textbf{REST / GraphQL} \\
\midrule
Perfil del usuario                &  1 & 1,8 kB & 0,08 kB &   1,8 / 0,1 kB \\
Publicaciones                     & 20 & 3,2 kB & 0,40 kB &  64,0 / 8,0 kB \\
Autor de cada publicación         & 20 & 1,8 kB & 0,08 kB &  36,0 / 1,6 kB \\
Comentarios                       & 60 & 1,1 kB & 0,12 kB &  66,0 / 7,2 kB \\
Autor de cada comentario          & 60 & 1,8 kB & 0,08 kB & 108,0 / 4,8 kB \\
\midrule
\textbf{Bytes transferidos $B$}   &    &        &         & \textbf{275,8 / 21,7 kB} \\
\textbf{Eficiencia de carga útil $\eta$} & & &            & \textbf{8 \% / 100 \%} \\
\textbf{Peticiones HTTP}          &    &        &         & \textbf{102 / 1} \\
\textbf{Viajes encadenados $n$}   &    &        &         &
  \textbf{17 (HTTP/1.1) o 4 (HTTP/2) / 1} \\
\bottomrule
\end{tabularx}
\end{table}

Tres cosas de ese cuadro merecen comentario.

La primera: mi $\eta$ para REST sale del 8 por ciento, mientras
\textcite{brito2019} midieron un 1 por ciento en bytes sobre sistemas reales.
Mi escenario es deliberadamente conservador y aun así la diferencia es de un
factor 13.

La segunda: las 102 peticiones no se convierten en 102 viajes. Sobre HTTP/1.1,
con seis conexiones simultáneas por servidor, son 17 rondas. Sobre HTTP/2, que
llegó en mayo de 2015 y multiplexa sobre una sola conexión
\parencite{rfc7540}, lo que queda son los cuatro niveles de dependencia del
grafo de datos. \textbf{HTTP/2 le quitó a GraphQL tres cuartas partes de su
ventaja el mismo año en que GraphQL se publicó.} Es el dato más incómodo de
todo este trabajo y volveré sobre él.

La tercera: $n=1$ es el mínimo teórico. GraphQL no puede mejorar más por ese
eje. Eso es lo que significa estar cerca de la asíntota.

\section{¿Cuánto tiempo transcurrió hasta que empezó a ser viable?}

La pregunta no tiene una respuesta sino cinco, porque \enquote{viable} significa
cosas distintas según quién pregunte. La Figura~\ref{fig:cronologia} las separa.

\begin{figure}[!htb]
\centering
\begin{tikzpicture}

% ================= Panel A: cronología y umbrales =================
\node[anchor=south west, font=\small\bfseries] at (-0.4,3.35)
     {Panel A. De prototipo interno a obligación contractual: trece años};

\draw[cGris!70, -{Stealth[length=5pt]}] (-0.15,0) -- (13.5,0);
\foreach \a in {2012,2014,2016,2018,2020,2022,2024,2026} {
  \draw[cGris!70] ({\tx{\a}},0) -- ({\tx{\a}},-0.10);
}
\node[below, font=\scriptsize, text=cGris] at ({\tx{2012}},-0.12) {2012};
\node[below right, font=\scriptsize, text=cGris]
     at ({\tx{2016}+0.08},-0.12) {2016};
\node[below, font=\scriptsize, text=cGris] at ({\tx{2020}},-0.12) {2020};
\node[below, font=\scriptsize, text=cGris] at ({\tx{2024}},-0.12) {2024};

\hito{2012}{1.05}{cGql}{\textbf{2012}\\en producción en iOS}{above}
\hito{2015}{2.15}{cGql}{\textbf{sep. 2015}\\especificación abierta}{above}
\hito{2016}{-1.05}{cRest}{\textbf{sep. 2016}\\GitHub, primera API
  pública}{below}
\hito{2018.5}{1.05}{cRef}{\textbf{2018 y 2019}\\Foundation y
  federación}{above}
\hito{2024.5}{2.15}{cVerde}{\textbf{2024 y 2025}\\Shopify obliga a
  GraphQL}{above}

\umbral{-2.30}{2012.4}{meses: viable \emph{dentro} de Facebook}{cGql}
\umbral{-2.78}{2015.7}{3 años: usable por terceros}{cGql}
\umbral{-3.26}{2016.7}{4 años: probada en público a escala}{cRest}
\umbral{-3.74}{2019.4}{7 años: existe andamiaje de empresa}{cRef}
\umbral{-4.22}{2024.8}{12 años: obligatoria en un dominio}{cVerde}

% ================= Panel B: curva S de difusión =================
\begin{scope}[yshift=-10.1cm]
\node[anchor=south west, font=\small\bfseries] at (-0.4,4.30)
     {Panel B. La curva S de difusión, que va muy por detrás de la de
      rendimiento};

\foreach \v in {25,50,75,100} {
  \draw[cGris!28] (0,{\dy{\v}}) -- (13.4,{\dy{\v}});
  \node[left, font=\scriptsize, text=cGris] at (-0.07,{\dy{\v}}) {\v};
}
\draw[cGris!70] (0,0) -- (13.4,0);
\draw[cGris!70] (0,0) -- (0,4.05);
\foreach \a in {2015,2018,2021,2024,2027,2030} {
  \draw[cGris!70] ({\dx{\a}},0) -- ({\dx{\a}},-0.09);
  \node[below, font=\scriptsize, text=cGris] at ({\dx{\a}},-0.09) {\a};
}
\node[left, font=\scriptsize, text=cGris, rotate=90, anchor=south]
     at (-0.62,2.0) {\% de empresas};

% REST, saturado
\draw[cRest, very thick, dash pattern=on 4pt off 2pt]
  (0,{\dy{93}}) -- (13.4,{\dy{93}});
\node[right, font=\scriptsize, text=cRest, align=left]
     at ({\dx{2015.2}},{\dy{97.5}})
     {REST: 93 \% en 2025, curva S terminada};

% GraphQL, logística ajustada a los dos puntos de Gartner
\draw[cGql, very thick, samples=80, domain=2014:2026, variable=\x]
  plot ({\dx{\x}},{\dy{100/(1+exp(-0.5565*(\x-2026.27)))}});
\draw[cGql, very thick, dashed, samples=50, domain=2026:2031, variable=\x]
  plot ({\dx{\x}},{\dy{100/(1+exp(-0.5565*(\x-2026.27)))}});

\fill[cGql] ({\dx{2021}},{\dy{9}}) circle (2.4pt);
\node[above left, font=\scriptsize, text=cGql, inner sep=2pt]
     at ({\dx{2021}},{\dy{9}}) {9 \% (2021)};
\fill[cGql] ({\dx{2024}},{\dy{28}}) circle (2.4pt);
\node[above left, font=\scriptsize, text=cGql, inner sep=2pt]
     at ({\dx{2024}},{\dy{28}}) {28 \% (2024)};
\draw[cGris, fill=white, thick] ({\dx{2025}},{\dy{33}}) circle (2.4pt);
\draw[cGris!70, thin] ({\dx{2019.6}},{\dy{56}}) -- ({\dx{2024.85}},{\dy{34.5}});
\node[right, font=\scriptsize, text=cGris, align=left, text width=3.55cm,
      inner sep=2pt]
     at ({\dx{2016.1}},{\dy{58}}) {33 \% de \emph{desarrolladores}\\(Postman:
      otra base, no comparable)};
\draw[cGql, fill=white, thick] ({\dx{2027}},{\dy{60}}) circle (2.4pt);
\node[above left, font=\scriptsize, text=cGql, align=right, inner sep=3pt]
     at ({\dx{2027}},{\dy{60}}) {previsión de Gartner:\\60 \% en 2027};

\draw[cAviso, dashed] ({\dx{2026.27}},0) -- ({\dx{2026.27}},{\dy{48}});
\node[right, font=\scriptsize, text=cAviso, inner sep=3pt]
     at ({\dx{2026.45}},{\dy{5}}) {inflexión de la difusión: 2026};

\node[right, font=\scriptsize\itshape, text=cGris, align=left]
     at ({\dx{2028.5}},{\dy{58}}) {extrapolación};
\end{scope}

\end{tikzpicture}
\caption{Panel A: hitos documentados y los cinco umbrales de viabilidad. Fechas
de \textcite{postman2024historia}, \textcite{graphql2015},
\textcite{github2016}, \textcite{techcrunch2018},
\textcite{apollo2024federacion}, \textcite{spec2025} y
\textcite{shopify2024rest}. Panel B: curva de difusión empresarial. Los puntos
llenos son las dos referencias de Gartner recogidas por \textcite{ibm2023gartner}
y \textcite{wundergraph2026}, menos del 10 por ciento en 2021 y menos del 30
por ciento en 2024; el punto blanco de 2027 es la previsión de Gartner de más
del 60 por ciento; el punto gris es Postman, que mide desarrolladores y no
empresas y por eso no es comparable \parencite{postman2025}. La logística está
ajustada a los dos puntos de Gartner y su asíntota se ha fijado en 100 por
falta de un dato que la acote, que es el supuesto más débil de la figura.}
\label{fig:cronologia}
\end{figure}

\begin{enumerate}
\item \textbf{Viable dentro de Facebook: meses.} Del prototipo de la primavera
  de 2012 a servir el muro de noticias en el verano de 2012
  \parencite{postman2024historia}. Fue tan rápido porque la tecnología
  precursora ya estaba: Facebook llevaba años operando TAO, una capa de datos
  con forma de grafo, y FQL, un lenguaje de consulta propio. GraphQL no inventó
  el grafo, inventó la forma de pedirlo.
\item \textbf{Usable por un tercero: tres años.} La especificación y la
  implementación de referencia se publican el 14 de septiembre de 2015
  \parencite{graphql2015}.
\item \textbf{Probada en público a escala: cuatro años.} GitHub publica su API
  GraphQL en septiembre de 2016 \parencite{github2016}. Fue la primera API
  GraphQL pública grande, y es la fecha que yo tomaría como el verdadero
  arranque de la curva S de difusión.
\item \textbf{Con andamiaje de empresa: siete años.} La gobernanza llega en
  noviembre de 2018 con la GraphQL Foundation bajo la Linux Foundation
  \parencite{techcrunch2018}, y la pieza que faltaba para organizaciones
  grandes, la federación, en mayo de 2019 \parencite{apollo2024federacion}.
  Netflix la adopta en 2020 y hoy opera un grafo federado con más de mil
  millones de peticiones diarias, más de diez mil tipos y campos y más de
  quinientos desarrolladores activos \parencite{netflix2020,infoq2022netflix}.
\item \textbf{Obligatoria en un dominio: doce años.} Shopify declara heredada
  su REST Admin API el 1 de octubre de 2024 y obliga a GraphQL en las
  aplicaciones públicas nuevas desde el 1 de abril de 2025
  \parencite{shopify2024rest,shopify2025graphql}. Ese es el momento en que
  GraphQL deja de ser una elección y pasa a ser una condición de contrato.
\end{enumerate}

Y si se cuenta la genealogía completa, la respuesta es otra: \textbf{veintiocho
años} desde \textcite{birrell1984} hasta la concepción de GraphQL en 2012,
con una generación nueva cada siete u ocho años.

Lo que me parece más útil de esta lista es el desfase entre el punto 2 y el
punto 5: \textbf{nueve años entre que la tecnología está terminada y que
alguien la impone.} La curva S de rendimiento y la curva S de difusión no van
juntas, y confundirlas es el error clásico de una hoja de ruta.

\section{La figura de mérito medida, y su sensibilidad}
\label{sec:sensibilidad}

Aquí está el núcleo del trabajo. Aplico la ecuación~\eqref{eq:fom} al escenario
del Cuadro~\ref{cuadro:escenario} en cinco entornos de red, desde el 3G de 2012
hasta dos servicios en el mismo centro de datos, y comparo cuatro
arquitecturas. Fijo $T_{s} = 50$ ms para todas, de modo que la comparación
aísle lo que la arquitectura de la interfaz sí controla.

\begin{figure}[!htb]
\centering
\begin{tikzpicture}

% ================= Panel A: descomposición en 3G =================
\node[anchor=south west, font=\small\bfseries] at (-3.5,3.72)
     {Panel A. Tiempo hasta la vista completa en la red de 2012
      (RTT 200 ms, 1,5 Mbit/s)};

\barraT{2.85}{3400}{1471}{50}{\textbf{REST} sobre HTTP/1.1\\102
  peticiones}{cRest}{4.921 ms}
\barraT{1.95}{800}{1471}{50}{\textbf{REST} sobre HTTP/2\\102
  peticiones}{cRest}{2.321 ms}
\barraT{1.05}{200}{116}{50}{\textbf{BFF a medida}\\1 endpoint por
  pantalla}{cSoap}{366 ms}
\barraT{0.15}{200}{116}{50}{\textbf{GraphQL}\\1 consulta}{cGql}{366 ms}

\draw[cAviso, thick, decorate,
      decoration={brace,amplitude=4pt,mirror}]
     (2.30,-0.20) -- (2.30,1.40);
\node[right, font=\scriptsize, text=cAviso, align=left, text width=8.6cm]
     at (2.44,0.60)
     {idénticos en esta figura de mérito: lo que separa a GraphQL de un
      \emph{backend for frontend} escrito a mano no es $T$, es $E$};

% leyenda
\fill[cRest] (0,-0.72) rectangle (0.42,-0.48);
\node[right, font=\scriptsize] at (0.47,-0.60) {viajes de ida y vuelta};
\fill[cRest!35] (3.95,-0.72) rectangle (4.37,-0.48);
\node[right, font=\scriptsize] at (4.42,-0.60) {transferencia de los bytes};
\fill[cGris!55] (8.45,-0.72) rectangle (8.87,-0.48);
\node[right, font=\scriptsize] at (8.92,-0.60) {trabajo del servidor (50 ms)};

% ================= Panel B: sensibilidad al RTT =================
\begin{scope}[yshift=-7.3cm]
\node[anchor=south west, font=\small\bfseries] at (-1.5,5.20)
     {Panel B. La sensibilidad: cuánto vale GraphQL frente a REST según la red};

\foreach \v in {2,4,6,8,10,12,14} {
  \draw[cGris!28] (0,{\py{\v}}) -- (12.3,{\py{\v}});
  \node[left, font=\scriptsize, text=cGris] at (-0.07,{\py{\v}})
       {\v$\times$};
}
\draw[cGris!70] (0,0) -- (12.3,0);
\draw[cGris!70] (0,0) -- (0,4.30);
\node[left, font=\scriptsize, text=cGris, rotate=90, anchor=south]
     at (-0.78,2.15) {ventaja de GraphQL};

\foreach \r/\t in {0.5/{0,5},1/1,5/5,20/20,50/50,200/200} {
  \draw[cGris!70] ({\px{\r}},0) -- ({\px{\r}},-0.09);
  \node[below, font=\scriptsize, text=cGris] at ({\px{\r}},-0.09) {\t};
}
\node[below, font=\scriptsize, text=cGris] at (6.1,-0.52)
     {tiempo de ida y vuelta de la red, en ms (escala logarítmica)};

% zona donde la ventaja deja de decidir
\fill[cVerde!8] (0,0) rectangle (12.3,{\py{1.5}});
\draw[cVerde!60, dashed] (0,{\py{1.5}}) -- (12.3,{\py{1.5}});
\node[right, font=\scriptsize\itshape, text=cVerde!60!black]
     at ({\px{2.2}},{\py{0.72}})
     {por debajo de 1,5$\times$ la figura de mérito ya no decide};

% curva frente a REST sobre HTTP/1.1
\draw[cRest, very thick]
  ({\px{0.5}},{\py{1.16}}) -- ({\px{5}},{\py{2.56}}) --
  ({\px{20}},{\py{5.75}}) -- ({\px{50}},{\py{9.29}}) --
  ({\px{200}},{\py{13.46}});
\foreach \r/\v in {0.5/1.16,5/2.56,20/5.75,50/9.29,200/13.46} {
  \fill[cRest] ({\px{\r}},{\py{\v}}) circle (2.4pt);
}
\node[right, font=\scriptsize, text=cRest, align=left]
     at ({\px{0.85}},{\py{10.4}}) {frente a REST\\sobre HTTP/1.1};

% curva frente a REST sobre HTTP/2
\draw[cGql, very thick]
  ({\px{0.5}},{\py{1.03}}) -- ({\px{5}},{\py{1.39}}) --
  ({\px{20}},{\py{2.12}}) -- ({\px{50}},{\py{3.31}}) --
  ({\px{200}},{\py{6.35}});
\foreach \r/\v in {0.5/1.03,5/1.39,20/2.12,50/3.31,200/6.35} {
  \fill[cGql] ({\px{\r}},{\py{\v}}) circle (2.4pt);
}
\node[right, font=\scriptsize, text=cGql, align=left]
     at ({\px{0.85}},{\py{6.3}}) {frente a REST\\sobre HTTP/2};

% escenarios de red rotulados
\foreach \r/\t in {200/{3G, 2012},50/{4G, 2016},20/{5G, 2022},
                   5/{fibra, 2026}} {
  \draw[cGris!45, dotted] ({\px{\r}},0) -- ({\px{\r}},4.42);
  \node[above, font=\scriptsize, text=cGris, align=center, inner sep=2pt]
       at ({\px{\r}},4.42) {\t};
}
\draw[cGris!45, dotted] ({\px{0.5}},0) -- ({\px{0.5}},4.42);
\node[above right, font=\scriptsize, text=cGris, align=left, inner sep=2pt]
     at ({\px{0.5}},4.42) {mismo centro de datos};
\end{scope}

\end{tikzpicture}
\caption{La ventaja de GraphQL no es una propiedad de GraphQL: es una función
del RTT. Panel A: descomposición de la ecuación~\eqref{eq:fom} para el
escenario del Cuadro~\ref{cuadro:escenario} en la red móvil de 2012. Panel B:
la misma comparación repetida en cinco entornos de red, cada uno con su RTT y
el ancho de banda típico de esa tecnología (1,5, 20, 100, 300 y 10.000
Mbit/s). Elaboración propia con $B = 275{,}8$ kB para REST, $B = 21{,}7$ kB
para GraphQL y $T_{s} = 50$ ms para ambos. Los cinco puntos son escenarios
discretos, no una función continua del RTT: la recta que los une solo ayuda a
leer la tendencia.}
\label{fig:sensibilidad}
\end{figure}

Los resultados están en la Figura~\ref{fig:sensibilidad} y son estos:

\begin{itemize}
\item En la red de 2012, 3G con 200 ms de ida y vuelta, GraphQL entrega la
  pantalla en 366 ms frente a los 4.921 ms de REST sobre HTTP/1.1. Es un
  \textbf{factor 13,5}, y explica sin ningún misterio por qué la tecnología
  nació en el equipo de móvil de Facebook y no en otro sitio.
\item Sobre HTTP/2 y todavía en 3G la ventaja cae a \textbf{6,4}. Sobre 4G, a
  \textbf{3,3}. Sobre 5G, a \textbf{2,1}. Sobre fibra doméstica, a
  \textbf{1,4}. Entre dos servicios del mismo centro de datos, a
  \textbf{1,03}.
\item El Panel A enseña además algo que no esperaba: \textbf{un endpoint
  compuesto a medida da exactamente el mismo tiempo que GraphQL}, 366 ms. En la
  figura de mérito principal son indistinguibles.
\end{itemize}

\subsection*{Lo que se aprende declarando bien la figura de mérito}

Los dos últimos puntos, juntos, cambian el argumento de venta de GraphQL.

Si la única figura de mérito fuera $T$, la conclusión sería que \textbf{GraphQL
solo se justifica cuando el consumidor está lejos}: aplicaciones móviles,
mercados con red pobre, clientes en otro continente. Y que entre servicios de
un mismo centro de datos no aporta nada, que es exactamente lo que ocurre en la
práctica y por qué ahí ganan REST y gRPC, que Google publicó en febrero de 2015
sobre HTTP/2 \parencite{grpc2015}.

Pero hay una segunda figura de mérito, $E$, el esfuerzo por consulta nueva, y
ahí sí se separan GraphQL y el endpoint a medida. \textcite{brito2020} hicieron
un experimento controlado con 22 estudiantes que implementaron ocho consultas
en las dos tecnologías: \textbf{la mediana fue de 9 minutos con REST y 6 con
GraphQL}, un 33 por ciento menos, y la ventaja crecía con la complejidad de la
consulta. El endpoint a medida, además, exige un despliegue del servidor por
cada pantalla nueva; GraphQL, ninguno.

Es la misma lección de la aviación: cuando una figura de mérito se satura, el
valor de la tecnología no desaparece, se muda a otra figura de mérito. La de
GraphQL se mudó del tiempo de respuesta al coste marginal de la vista número
cincuenta.

Conviene añadir que la literatura empírica no es unánime, y que eso encaja con
el resultado de la sensibilidad en lugar de contradecirlo.
\textcite{seabra2019} compararon servicios antes y después de migrarlos a
GraphQL: hubo mejora en dos de cada tres casos, pero por encima de unas tres
mil peticiones por prueba el servicio en GraphQL rindió \emph{por debajo} de su
equivalente en REST. Es exactamente lo que predice la
ecuación~\eqref{eq:fom} cuando el cuello de botella deja de estar en la red y
pasa a estar en el servidor, y es el asunto de la sección siguiente.

\section{Restricciones activas del diseño}
\label{sec:restricciones}

La pregunta pide las restricciones activas que limitan el rendimiento del
sistema. En el lenguaje del módulo, una restricción es activa cuando está
ligada en el óptimo: relajarla mejora la figura de mérito, y relajar cualquier
otra no hace nada. Con esa definición, y a la vista de la
Figura~\ref{fig:sensibilidad}, la lista de GraphQL ha cambiado de orden en los
últimos diez años.

\begin{table}[t]
\small
\caption{Las cinco restricciones activas del diseño de GraphQL, ordenadas por
cuánto limitan hoy la figura de mérito. Las tres primeras comparten mitigación,
y esa mitigación tiene un coste común.}
\label{cuadro:restricciones}
\begin{tabularx}{\textwidth}{@{}p{3.3cm}XXp{3.3cm}@{}}
\toprule
\textbf{Restricción activa} & \textbf{De dónde sale} &
\textbf{Qué límite impone} & \textbf{Mitigación y su coste} \\
\midrule
\textbf{1. Resolución campo a campo}
& La especificación resuelve cada campo por separado; una lista de $n$
  elementos con un campo que consulta el almacén genera $n+1$ accesos
& Gobierna $T_{s}$, que en redes rápidas es ya el 98 \% de $T$
& DataLoader agrupa por ciclo del bucle de eventos, o se compila el árbol a
  SQL. Coste: es responsabilidad del programador, no de la plataforma \\
\addlinespace[2pt]
\textbf{2. Espacio de consulta no acotado}
& El cliente compone; el coste no se puede leer de la URL
& La misma ruta puede costar 1 unidad o un millón. Escape halló problemas de
  consumo de recursos no restringido en el 69 \% de 160 endpoints públicos
  analizados \parencite{escape2024}
& Coste calculado por consulta: Shopify da 1.000 puntos por ventana y repone
  entre 50 y 500 por segundo \parencite{shopifylimites}. Coste: la capacidad
  pasa a ser un problema de esquema \\
\addlinespace[2pt]
\textbf{3. Sin caché HTTP intermedia}
& POST a un único punto de entrada: ni CDN, ni proxy, ni caché del navegador
& Pierde el acierto de borde, que en REST es gratis y llega al 70 u 80 \%
& Consultas persistidas con GET y hash. Coste: vuelve a atar el cliente al
  despliegue del servidor \\
\addlinespace[2pt]
\textbf{4. Errores fuera de HTTP}
& Se responde 200 con un array \texttt{errors}
& Cortafuegos, balanceadores y observabilidad por código de estado dejan de
  ver los fallos
& Tipos de resultado en el esquema. Coste: disciplina, no tecnología \\
\addlinespace[2pt]
\textbf{5. Gobernanza del esquema}
& Un grafo único compartido por muchos equipos
& Fija el tamaño mínimo de organización a partir del cual compensa
& Federación, registro de esquemas y comprobación de composición. Coste:
  una plataforma y un equipo que la sostenga \\
\bottomrule
\end{tabularx}
\end{table}

\subsection*{La restricción dominante se mudó de la red al servidor}

Esta es, para mí, la conclusión más importante del trabajo. En 2012 el término
que gobernaba la ecuación~\eqref{eq:fom} era $n \cdot \mathrm{RTT}$: 3.400 de
los 4.921 ms, el 69 por ciento. GraphQL atacó ese término y lo llevó a su
mínimo teórico, $n=1$.

Hoy, entre dos servicios del mismo centro de datos, ese mismo término vale
0,5 ms de un total de 50,5. El 99 por ciento del tiempo es $T_{s}$.
\textbf{GraphQL agotó el eje que optimizaba, y la restricción activa pasó a ser
precisamente su propio modelo de ejecución.}

Y ese modelo de ejecución tiene un problema estructural, no incidental: la
especificación resuelve cada campo de forma independiente, de modo que el
resolvedor de un campo dentro de una lista de $n$ elementos se ejecuta $n$
veces. El Panel A de la Figura~\ref{fig:restricciones} lo dibuja. La solución
canónica, DataLoader, agrupa las claves pedidas dentro de un mismo ciclo del
bucle de eventos y las resuelve en un solo acceso \parencite{dataloader}. Es
eficaz, pero conviene ver lo que significa: \textbf{el rendimiento de una API
GraphQL no depende de la tecnología sino de si el equipo puso el DataLoader}.
En REST ese error no se puede cometer, porque el endpoint se escribe entero a
mano.

\begin{figure}[!htb]
\centering
\begin{tikzpicture}

% ================= Panel A: el N+1 de resolvedores =================
\node[anchor=south west, font=\small\bfseries] at (-0.9,5.02)
     {Panel A. La restricción activa de hoy};
\node[anchor=south west, font=\scriptsize\itshape, text=cGris] at (-0.9,4.70)
     {accesos al almacén de datos};

\foreach \v in {25,50,75,100} {
  \draw[cGris!28] (0,{\ny{\v}}) -- (6.45,{\ny{\v}});
  \node[left, font=\scriptsize, text=cGris] at (-0.07,{\ny{\v}}) {\v};
}
\draw[cGris!70] (0,0) -- (6.45,0);
\draw[cGris!70] (0,0) -- (0,4.55);
\foreach \a in {20,40,60,80,100} {
  \draw[cGris!70] ({\nx{\a}},0) -- ({\nx{\a}},-0.09);
  \node[below, font=\scriptsize, text=cGris] at ({\nx{\a}},-0.09) {\a};
}
\node[below, font=\scriptsize, text=cGris] at (3.1,-0.48)
     {tamaño de la lista, $n$};

\draw[cAviso, very thick] ({\nx{0}},{\ny{1}}) -- ({\nx{100}},{\ny{101}});
\node[left, font=\scriptsize, text=cAviso, align=right]
     at ({\nx{88}},{\ny{92}}) {resolvedor ingenuo:\\$n+1$ accesos};

\draw[cVerde, very thick] ({\nx{0}},{\ny{2}}) -- ({\nx{100}},{\ny{2}});
\node[right, font=\scriptsize, text=cVerde, align=left]
     at ({\nx{37}},{\ny{18}}) {con DataLoader: 2 accesos,\\constante en $n$};

\draw[cGris!80, dashed] ({\nx{20}},0) -- ({\nx{20}},{\ny{21}});
\fill[cAviso] ({\nx{20}},{\ny{21}}) circle (2.2pt);
\fill[cVerde] ({\nx{20}},{\ny{2}}) circle (2.2pt);
\draw[cGris!80, thin] ({\nx{28}},{\ny{74}}) -- ({\nx{20.7}},{\ny{25}});
\node[right, font=\scriptsize, text=cGris, align=left, inner sep=3pt]
     at ({\nx{5}},{\ny{84}}) {mi escenario:\\21 accesos\\frente a 2};

% ================= Panel B: frente de Pareto =================
\begin{scope}[xshift=9.3cm]
\node[anchor=south west, font=\small\bfseries] at (-0.9,5.02)
     {Panel B. Por qué las tres conviven};
\node[anchor=south west, font=\scriptsize\itshape, text=cGris] at (-0.9,4.70)
     {coste de servidor y de gobernanza (índice propio)};

\foreach \v in {25,50,75,100} {
  \draw[cGris!28] (0,{\fy{\v}}) -- (6.45,{\fy{\v}});
  \node[left, font=\scriptsize, text=cGris] at (-0.07,{\fy{\v}}) {\v};
}
\draw[cGris!70] (0,0) -- (6.45,0);
\draw[cGris!70] (0,0) -- (0,4.80);
\foreach \a in {25,50,75,100} {
  \draw[cGris!70] ({\fx{\a}},0) -- ({\fx{\a}},-0.09);
  \node[below, font=\scriptsize, text=cGris] at ({\fx{\a}},-0.09) {\a};
}
\node[below, font=\scriptsize, text=cGris] at (3.1,-0.48)
     {eficiencia de carga útil $\eta$, en \%};

% región dominada
\fill[cAviso!7] ({\fx{8}},{\fy{20}}) -- ({\fx{35}},{\fy{20}}) --
                ({\fx{35}},{\fy{25}}) -- ({\fx{98}},{\fy{25}}) --
                ({\fx{98}},{\fy{45}}) -- ({\fx{8}},{\fy{45}}) -- cycle;

% frente
\draw[cVerde, very thick, dash pattern=on 5pt off 2pt]
  ({\fx{8}},{\fy{20}}) -- ({\fx{35}},{\fy{25}}) -- ({\fx{98}},{\fy{45}});
\node[below right, font=\scriptsize\itshape, text=cVerde, align=left]
     at ({\fx{60}},{\fy{30}}) {frente de Pareto};

\punto{4}{75}{cSoap}{SOAP y WS-$\ast$}{above right}
\punto{8}{20}{cRest}{REST y JSON}{below right}
\punto{35}{25}{cRpc}{gRPC\\\emph{solo servidor}\\\emph{a servidor}}{above left}
\punto{95}{65}{cGris}{BFF a medida}{above left}
\punto{98}{45}{cGql}{GraphQL}{above}

\draw[cAviso!65, -{Stealth[length=3.5pt]}, thin]
     ({\fx{44}},{\fy{52}}) -- ({\fx{47}},{\fy{40}});
\node[font=\scriptsize\itshape, text=cAviso, align=center]
     at ({\fx{44}},{\fy{58}}) {zona dominada};
\end{scope}

\end{tikzpicture}
\caption{Panel A: el problema $N+1$ de resolvedores, que es la restricción
activa dominante de GraphQL una vez que la red deja de serlo. La curva ingenua
es el comportamiento que describe la propia documentación de referencia; la
constante es lo que consigue DataLoader agrupando por ciclo del bucle de
eventos \parencite{dataloader}. Panel B: las cinco arquitecturas situadas
frente a las dos figuras de mérito que de verdad compiten. Las posiciones en el
eje vertical son un índice ordinal de elaboración propia, no una medida: lo que
esta figura afirma es la \emph{forma} del frente, no las coordenadas. SOAP y el
\emph{backend for frontend} escrito a mano están dominados; REST, gRPC y
GraphQL están los tres sobre el frente, cada uno en un tramo distinto, y por
eso conviven.}
\label{fig:restricciones}
\end{figure}

\subsection*{Las tres primeras restricciones comparten medicina, y la medicina
tiene un precio}

La consulta persistida resuelve a la vez la 2 y la 3, y ayuda con la 1: el
cliente envía un hash en lugar de la consulta, el servidor la busca en un
registro, y como es un GET conocido y de coste acotado vuelven a funcionar la
caché de borde y la limitación de tasa.

Es la mejor mitigación que existe y por eso hay que ver lo que cuesta. Con
consultas persistidas, el cliente ya no puede pedir cualquier cosa: solo puede
pedir lo que está registrado, y registrar es un despliegue coordinado.
\textbf{Es decir, la mitigación devuelve buena parte del acoplamiento entre
cliente y servidor que GraphQL vino a eliminar.} No es una contradicción, es un
frente de Pareto: se compra rendimiento y seguridad con flexibilidad, y cada
organización elige su punto.

El Panel B de la Figura~\ref{fig:restricciones} sitúa las cinco arquitecturas
sobre ese frente. El resultado es el que explica el estado del mercado mejor
que cualquier argumento: \textbf{REST, gRPC y GraphQL están los tres sobre el
frente de Pareto.} Ninguno domina a los otros. SOAP y el endpoint a medida sí
están dominados. Por eso REST sigue en el 93 por ciento y GraphQL en el 33
\parencite{postman2025}, y por eso la pregunta \enquote{¿GraphQL sustituirá a
REST?} está mal planteada: no se sustituye a un punto del frente, se elige otro
punto.

\section{Dónde está GraphQL hoy y qué haría yo con esto}

Poniendo las tres curvas juntas, la lectura es esta:

\begin{enumerate}
\item \textbf{La curva S de rendimiento está prácticamente terminada.} $n=1$ es
  el mínimo teórico y $\eta \approx 1$ es el máximo teórico. Por esos dos ejes
  no queda nada que ganar. Lo que se publica hoy en el ecosistema no mejora la
  figura de mérito original: mejora el $T_{s}$ y la gobernanza.
\item \textbf{La curva S de difusión está justo en su punto de inflexión}, en
  torno a 2026 según el ajuste del Panel B de la
  Figura~\ref{fig:cronologia}. La adopción sigue creciendo nueve años después
  de que la tecnología estuviera terminada, que es el desfase normal entre las
  dos curvas.
\item \textbf{La ventana de sensibilidad se ha estrechado por el lado de la
  red.} HTTP/2 en 2015 y luego 4G y 5G se llevaron un orden de magnitud de la
  ventaja. La ventana sigue abierta donde el RTT es alto: móvil en mercados con
  red pobre, clientes intercontinentales, dispositivos de baja potencia.
\end{enumerate}

Traducido a una decisión, que es para lo que sirve esto: \textbf{yo no
justificaría hoy una migración a GraphQL con el argumento del tiempo de
respuesta}, salvo que el consumidor esté medido y esté lejos. La
Figura~\ref{fig:sensibilidad} dice que ese argumento vale un factor 1,4 sobre
fibra. Los argumentos que sí aguantan la aritmética son el 33 por ciento de
esfuerzo por consulta de \textcite{brito2020}, la eliminación del despliegue
por pantalla nueva y la gobernanza de un esquema único cuando hay muchos
equipos y muchos clientes, que es el caso de Netflix y de Shopify. Y ninguno de
los tres es un argumento de rendimiento.

\section{Limitaciones}

\begin{enumerate}
\item \textbf{El escenario del Cuadro~\ref{cuadro:escenario} es sintético.} Los
  tamaños de objeto son mi estimación de un caso típico, no una medición. Elegí
  valores conservadores, y aun así la $\eta$ que sale para REST, el 8 por
  ciento, es ocho veces mejor que la mediana medida por
  \textcite{brito2019} sobre sistemas reales. Con sus números todas mis
  ventajas serían mayores.
\item \textbf{Las cinco curvas de la Figura~\ref{fig:curvass} son logísticas
  ajustadas a mano} sobre fechas documentadas de aparición y saturación, no
  series observadas. Sirven para leer la forma y el relevo, no un valor.
\item \textbf{Fijar $T_{s} = 50$ ms para las cuatro arquitecturas es una
  simplificación fuerte y favorable a REST en un sentido y a GraphQL en otro.}
  Un GraphQL sin DataLoader tendría un $T_{s}$ mucho peor, que es precisamente
  la restricción activa de la sección~\ref{sec:restricciones}, y un REST con
  102 peticiones consume mucha más CPU de servidor por vista.
\item \textbf{Los cinco escenarios de red del Panel B son puntos discretos},
  cada uno con su pareja de RTT y ancho de banda. La línea que los une no es
  una función continua y no debe leerse como tal.
\item \textbf{La asíntota de la curva de difusión está fijada en 100 por
  falta de un dato que la acote.} Es el supuesto más débil de todo el trabajo.
  Si el techo real fuera el 70 por ciento, el punto de inflexión se adelantaría
  y la previsión de Gartner para 2027 quedaría muy cerca de la saturación.
\item \textbf{Los dos puntos de Gartner y el de Postman no son comparables}: el
  primero mide porcentaje de empresas con GraphQL en producción y el segundo
  porcentaje de desarrolladores encuestados. Van juntos en el Panel B para dar
  contexto, no para ajustar la curva, y por eso el de Postman está dibujado en
  gris.
\item \textbf{El eje vertical del frente de Pareto es ordinal y es mío.} Está
  construido con mi criterio sobre coste de servidor y de gobernanza. La figura
  afirma la forma del frente y qué puntos están dominados, no las coordenadas.
\item \textbf{No he medido nada yo.} Todo lo cuantitativo viene de la
  literatura citada o de aplicar la ecuación~\eqref{eq:fom} a un escenario
  declarado. Un banco de pruebas propio sobre una API real sería el siguiente
  paso y probablemente cambiaría los detalles, aunque no el orden de magnitud
  ni la forma de la curva de sensibilidad, que es lo que sostiene la
  conclusión.
\end{enumerate}

\printbibliography

\end{document}
